\documentclass{article}
\usepackage{graphicx}
\usepackage{booktabs}
\usepackage{amsmath}
\usepackage{hyperref}
\usepackage{multirow}

\title{RetrievalBench: Cross-Domain Benchmarking of Retrieval Techniques for Long Structured Documents}
\author{[Intern Name], Spurthi Setty \\
Anote AI}
\date{}

\begin{document}

\maketitle

\begin{abstract}
\noindent
Retrieval-Augmented Generation (RAG) systems depend critically on retrieval quality, yet existing evaluations are domain-specific and provide little guidance on whether winning techniques generalize. We extend our prior work on financial document retrieval \cite{setty2024rag} to a cross-domain benchmark spanning finance (FinanceBench), legal (CUAD), medical (PubMedQA), and technical (TechQA) documents. We evaluate 24 configurations from a six-level ablation of retrieval techniques: chunking strategy, query expansion, re-ranking, embedding type, metadata annotation, and hybrid search. Our results show that the optimal retrieval pipeline varies substantially by domain, and that the financial-domain winner transfers to legal documents but not to medical or technical corpora. We release RetrievalBench, including all datasets, code, and a cross-domain retrieval technique leaderboard.
\end{abstract}

\section{Introduction}
% RAG retrieval is domain-dependent; most work evaluates one domain
% Our prior work (arXiv:2404.07221) achieved strong results on finance
% This paper: does the winner generalize? Cross-domain benchmark

\section{Related Work}
\subsection{RAG Retrieval Techniques}
\subsection{Domain-Specific Document QA}
\subsection{Cross-Domain Generalization in NLP}

\section{RetrievalBench}
\subsection{Domains and Datasets}
\begin{table}[h!]
\centering
\caption{RetrievalBench Domains}
\begin{tabular}{llll}
\toprule
\textbf{Domain} & \textbf{Dataset} & \textbf{Doc Type} & \textbf{\# Questions} \\
\midrule
Finance & FinanceBench & 10-K, 10-Q, 8-K & 150 \\
Legal & CUAD & Contracts & [N] \\
Medical & PubMedQA & Research Articles & [N] \\
Technical & TechQA & Tech Manuals & [N] \\
\bottomrule
\end{tabular}
\end{table}

\subsection{Six-Level Retrieval Ablation}
% Level 1: Chunking (fixed / recursive / semantic / element-based)
% Level 2: Query Expansion (none / HyDE / multi-query)
% Level 3: Re-ranking (none / cross-encoder / Cohere)
% Level 4: Embedding (generic / domain fine-tuned)
% Level 5: Metadata annotation (none / auto-annotated)
% Level 6: Hybrid search (dense-only / dense+sparse BM25)

\section{Experimental Setup}
\subsection{Evaluation Metrics}
% Page-level retrieval accuracy, LLM evaluation score, BLEU, ROUGE-L
\subsection{Models}
% Retriever: various; Generator: GPT-4o

\section{Results}
% Table: 24-condition ablation results per domain
% Figure: Cross-domain performance heatmap
% Figure: Comparison with arXiv:2404.07221 baseline (finance)

\subsection{Best Configuration Per Domain}
\subsection{Cross-Domain Transfer Analysis}
\subsection{Ablation: Which Technique Matters Most?}

\section{Analysis}
\subsection{Why Finance Techniques Transfer to Legal but Not Medical}
\subsection{Document Structure as a Predictor of Technique Effectiveness}

\section{Conclusion}

\section*{References}
\small
\begin{enumerate}
    \item Setty et al. Improving Retrieval for RAG based Question Answering Models on Financial Documents, 2024. arXiv:2404.07221.
    \item Gao et al. Retrieval-Augmented Generation for Large Language Models: A Survey, 2024. arXiv:2312.10997.
    \item Patronus AI. FinanceBench, 2023.
    \item Hendrycks et al. CUAD: An Expert-Annotated NLP Dataset for Legal Contract Review, 2021.
    \item Jin et al. PubMedQA: A Biomedical Research Question Answering Dataset, 2019.
    % Add more references
\end{enumerate}

\end{document}
